%With the breakthrough of Large Language Models (LLMs) in question answering, code generation, and multimodal understanding, Retrieval-Augmented Generation (RAG) has become the mainstream paradigm for improving model factual consistency and knowledge coverage. RAG effectively alleviates the problems of "knowledge timeliness" and "hallucination generation" in language models by combining external knowledge base retrieval with generative models. However, most mainstream RAG systems rely on LLMs with tens of billions of parameters and expensive embedded retrieval engines. This introduces two critical challenges that significantly limit practical deployment.

%On the one hand, the LLM-RAG framework exhibits significant resource-performance contradictions when dealing with long documents versus dynamic knowledge. Due to the high dimensionality of embedding vectors and the large redundancy of retrieval databases, the system often faces delay accumulation and memory bottlenecks during the recall stage. On the other hand, the semantic breakage and context loss caused by the text chunking strategy have not been effectively solved, resulting in incoherent logic or incomplete evidence when generating answers. In addition, resource constraints on edge devices or in enterprise private deployments further limit the scalability and maintainability of traditional RAG.

With the breakthrough of Large Language Models (LLMs) in question answering, code generation, and multimodal understanding~\citep{gao2023retrieval,wang2025think,wang2026efficient,li2026arise}, Retrieval-Augmented Generation (RAG)~\citep{lewis2020retrieval} has become the mainstream paradigm for improving model factual consistency and knowledge coverage.
RAG effectively alleviates the problems of knowledge timeliness and hallucination generation in language models by combining external knowledge base retrieval with generative models~\citep{wang2023augmenting,zhang2026tdarc}. 
Specifically, RAG systems consist of three core steps, namely retrieval, chunking, and generation.
First, the user provides an input query to the retriever, which derives retrieval keywords for sparse retrieval or rewrites the query for dense retrieval.
The retriever then extracts uniform-length chunks of information that are highly similar to the retrieval query.
The retrieved evidence is then aggregated and sent to the LLM as a generator to produce an answer to the user's query.


The standard RAG mechanisms, however, are prone to information insufficiency in the retrieved evidence.
The fixed retrieval channel (e.g., dense~\cite{borgeaud2022improving, izacard2022few, lewis2020retrieval} or sparse~\cite{ram2023context, huly2024old}) commonly leveraged in the contemporary RAG systems tends to ignore the user's requirement~\cite{wang2024evaluating}.
Such single-channel retrieval may collect descriptive evidence for semantically complex queries or retrieve explanatory evidence for purely descriptive questions, leading to evidence–intent mismatches.
Moreover, mechanical chunking (e.g., boundary-based chunking~\cite{finardi2024chronicles} or structure-aware chunking~\footnote{https://python.langchain.com}~\cite{zhang2026text}) may break the content or ignore the semantic connections, posing a risk to information fragmentation~\citep{izacard2021leveraging}.
Recent studies on efficient neural modeling and representation learning further suggest that rigid processing pipelines can hinder the preservation of temporal and structural dependencies in sequential data, motivating more adaptive and context-aware retrieval and encoding strategies~\citep{8368077,9540752}.

\begin{comment}


Most contemporary RAG systems rely on \say{Large} LM, which demands high computational resources.
Models' immense parameters limit the performance when handling long documents or dynamic knowledge sources. 
%and expensive embedded retrieval engines,
Further, the high dimensionality of embedding vectors and redundancy of retrieval databases lead to delay accumulation and memory bottlenecks during retrieval.
%This computational burden involves high inference latency, substantial memory footprints, and significant energy consumption, making deployment on edge devices or in enterprise private settings particularly challenging.
Recent lightweight RAG approaches such as MiniRAG~\citep{fan2025minirag} and LightRAG~\citep{guo2024lightrag} attempt to reduce computational requirements through data structure optimization and efficient routing. 
However, these methods primarily focus on architectural simplifications rather than embedding semantic reasoning throughout the entire pipeline.

\end{comment}


%The high dimensionality of embedding vectors and redundancy of retrieval databases lead to delay accumulation and memory bottlenecks during retrieval.
%This computational burden involves high inference latency, substantial memory footprints, and significant energy consumption, making deployment on edge devices or in enterprise private settings particularly challenging.
%Recent lightweight RAG approaches such as MiniRAG~\citep{fan2025minirag} and LightRAG~\citep{guo2024lightrag} attempt to reduce computational requirements through data structure optimization and efficient routing. However, these methods primarily focus on architectural simplifications rather than embedding semantic reasoning throughout the entire pipeline.


\begin{comment}

RAG retrieves chunks of relevant information from the external knowledge base and and concatenate them under the same chunk length to simplify the processing in the LM.
%The second critical challenge concerns evidence integrity and context preservation.
%Text chunking, a necessary step for indexing large documents, inevitably introduces semantic fragmentation. When documents are split into fixed-length chunks, retrieved passages often become truncated mid-sentence or disconnected from their surrounding context, depriving the model of critical information for coherent reasoning.
To achieve the same chunk length, chunks are truncated automatically in the conventional RAG systems.
This truncation introduces semantic fragmentation, jeopardizing the overall contextual information from the original full-length document~\cite{gunther2024late}.
%While recent work has explored improved chunking strategies at the indexing stage, such as semantic-aware chunking and hierarchical segmentation, few systems provide post-retrieval mechanisms to detect or repair evidence truncation. Additionally, even when individual pieces of evidence are coherent, assembling them into a comprehensive, consistent answer that covers all dimensions of a complex query remains problematic. Query enhancement techniques like DMQR-RAG~\citep{li2024dmqr} and HyDE~\citep{gao2022precise} improve retrieval diversity through multi-perspective search, but without systematic coverage verification, knowledge gaps often persist undetected. Furthermore, advanced RAG systems employing agentic mechanisms and closed-loop reasoning, such as Self-RAG~\citep{asai2024self} and CRAG~\citep{yan2024corrective}, improve autonomy while continuing to rely on large model orchestration, maintaining prohibitive computational costs.
%While 
%Recent work explores chunking strategies at the indexing stage, and a few of them provide post-retrieval mechanisms to detect or repair evidence truncation. 
%For example, 
To alleviate this issue, recent work introduce query enhancement techniques like DMQR-RAG~\citep{li2024dmqr} and HyDE~\citep{gao2022precise} to improve retrieval diversity.
Some more advanced systems, such as Self-RAG~\citep{asai2024self} and CRAG~\citep{yan2024corrective}, introduce autonomous mechanisms to improve chunking at the indexing stage.
Nevertheless, these approaches either rely on LLM or exclude the post-retrieval repair mechanism.
%, but without systematic coverage verification, knowledge gaps persist. Advanced RAG systems such as Self-RAG~\citep{asai2024self} and CRAG~\citep{yan2024corrective} introduce autonomous mechanisms, yet continue to rely on large model orchestration, maintaining prohibitive computational costs.




\end{comment}

%%%%%%%%%%%%%%%%%%%%%%%%
%To solve the above challenges, this paper proposes a new retrieval-augmented generation framework based on Small Language Model (SLM) - SLM-RAG. Unlike traditional LLM-RAG, which relies on ultra-large parameter scale, SLM-RAG builds a multi-layer semantic inference pipeline with a lightweight model as the core, and completes the whole process of autonomous decision-making through a set of collaborative SLM submodules. The core idea of this framework is to introduce the Hypothetical Query Reasoning mechanism, so that the model can actively generate multi-perspective semantic subqueries before retrieval, thereby improving the recall breadth and semantic coverage. After searching, the forensic stitching mechanism is automatically detected and repaired to ensure the integrity of the evidence context. At the same time, the built-in coverage audit module of the system can align and analyze the entities, relationships and time dimensions between problems and evidence, automatically trigger gap reretrieval and partial completion, and build a high-precision and self-correcting knowledge retrieval closed loop.

To address these challenges, we propose InSemRAG, a RAG framework that addresses the problem of information insufficiency in the retrieved evidence. 
%
%We are motivated by the excellent instruction following capability at computational cost brought by SLMs~\cite{wang2025comprehensive,lee2024mentor}.
%This enables our framework to perform simple semantic decomposition that enriches the retrieval skills of our RAG.
%
%The core insight is that SLMs can effectively perform sophisticated semantic reasoning at each pipeline stage when equipped with principled architectural mechanisms and structured task instructions. 
%SLM-RAG builds a multi-layer semantic inference pipeline with a lightweight model as the core, and completes the whole process of autonomous decision-making through a set of collaborative SLM submodules. 
%Unlike lightweight systems that optimize individual components or agentic systems that require large models for orchestration, SLM-RAG achieves both computational efficiency and autonomous reasoning by embedding semantic control throughout the entire pipeline.
%
%With an SLM as the main actor performing the retrieval task, we introduce two core modules that construct our framework.
InSemRAG consists of two modules, an intention-aware retriever (IAR) and semantics-preserving chunking (SPC).
Given an input query from the user, we first perform dual-view rewriting that simultaneously decomposes the query into keywords for sparse retrieval and expands it into a semantically meaningful phrase for dense retrieval. 
Utilizing the rewritten query, we perform dynamic hybrid retrieval, where we dynamically modulate the retrieval channel weight depending on the characteristics of the query.
Instead of passing the retrieved evidence directly to the generator, we perform detection and reparation on the semantically damaged chunk.
Specifically, we repair the damaged chunk by aggregating with the nearest neighbor chunks from the source documents and rewriting it to meet the chunk length requirement while maintaining the semantic information.
Finally, we assess whether the final evidence contains all query-required information elements.
Otherwise, the retrieve-and-check procedure is repeated until information completeness is satisfied.
%First, we apply \textbf{query-adaptive retrieval (QAR)} mechanisms that effectively adjust the type of retrieval based on what the query requires.
%Specifically, we enrich the user query with dense expansion and sparse decomposition.
%With such a multi-view improved query, we perform dual-retrieval whose channel contribution is controlled by dynamic weights.
%SLM-RAG introduces three key innovations to solve the identified problems. First, Hypothetical Query Reasoning generates multi-perspective semantic subqueries before retrieval, establishing a retrieval hypothesis space that expands coverage across different information dimensions and semantic angles. 
%This proactive approach improves recall breadth without requiring large model involvement. 
%We complement this adaptive retrieval with our second innovation, which is \textbf{semantic-preserving stitching (SPS)}.
%The main objective of this module is to check whether the chunks of evidence being passed to the LLM preserve the integrity contained in the source document.
%Specifically, we first detect the presence of any semantic damage in each chunk and repair it through stitching if it exists.
%Thus, we construct this module to alleviate the problem of blind chunking in the conventional RAG systems that result into unintentional information cutoff, jeopardizing the generation quality of the LLM.
%We design the whole retrieval process as an iterative process that continuously checks the quality of retrieval knowledge before being passed to the generator.
%Second, Forensic Evidence Completion automatically detects semantic truncation in retrieved passages post-retrieval and restores context through adjacent segment retrieval and extractive compression. This mechanism ensures evidence coherence without introducing hallucinations, directly addressing the evidence integrity problem. Third, Systematic Coverage Auditing through entity-relationship-temporal triple alignment automatically detects knowledge gaps between queries and retrieved evidence, triggering targeted re-retrieval to fill missing information. This creates a self-healing closed loop where gaps automatically trigger refinement, achieving autonomous reasoning without large model orchestration.

%While maintaining overall lightweight, SLM-RAG demonstrates excellent performance and scalability. Experimental results show that the framework significantly improves the retrieval accuracy and answer consistency on multiple open domain Q&A and enterprise knowledge base datasets. Compared to mainstream LLM-RAG, SLM-RAG reduces response latency by approximately 45\% with approximately 90\% reduction in computational overhead, and significantly improves evidence coverage and answer attributability.

To accommodate the complexity of the iterative retrieve-and-check mechanism in InSemRAG, we leverage small language models (SLMs) as the machines for retrieval and semantic checking.
We are motivated by the excellent instruction following capability at minimum computational cost brought by SLMs~\cite{wang2025comprehensive,lee2024mentor}.
%Furthermore, prior works have also investigated the effectiveness of SLMs for RAG~\cite{fan2025minirag, richard2025slm}.
This utilization circumvents the need for large instruction-following models utilized in recent RAG systems~\cite{shen2024retrieval}.

We summarize our contributions as follows:
%Through the incorporation of SLM in InSemRAG, we computational overhead by approximately 90\% compared to LLM-RAG systems while improving evidence quality, interpretability, and response latency by 45\%.
%The main contributions of this paper can be summarized as follows:
\begin{itemize}
    \item To the best of our knowledge, we are the first to introduce a semantic integrity-based iterative dynamic retrieve-and-check mechanism in RAG.
    
    \item We propose InSemRAG, an efficient intent-aware retrieval and semantics-preserving chunking mechanism for RAG to retrieve evidence that not only considers the user's requirement, but also maintains the evidence's semantic integrity.
    
    %\item We leverage SLMs as the retriever to alleviate the computational complexity of iterative retrieve-and-check mechanism of InSemRAG.
    %to perform retrieval on the RAG system, alleviating the computational bottleneck on conventional RAG systems that utilize LLM.
    %propose A novel paradigm of retrieval-augmented generation with small language models as the core, constructing a lightweight, autonomous, and explainable reasoning closed loop through multi-layer semantic inference.
    %\item We propose query-adaptive retrieval and semantic-preserving stitching that not only steer the retrieval behavior to match with the query's need but also preserve the semantic integrity of the retrieved evidence.
    %Introduction of Hypothetical Query Reasoning and Coverage Auditing mechanisms that systematically solve structural shortcomings of traditional RAG in terms of recall coverage and evidence integrity.
    %\item A forensic evidence completion strategy designed to effectively alleviate semantic rupture caused by text chunking, detecting and restoring truncated passages without hallucination.
    \item Through a comprehensive experimental validation, we demonstrate the effectiveness of InSemRAG in both performance and efficiency dimensions. Specifically, we yield up to 2.65-point improvement in F1 on evidence-sensitive reasoning tasks with 4.32× lower latency than prior state-of-the-art methods.
    
    %efficiency and accuracy advantages of our method across multiple datasets, providing a feasible path for intelligent question answering and knowledge retrieval in low-resource environments.
\end{itemize}

%In summary, our method demonstrates the potential capabilities of SLMs in complex retrieval and generation tasks, providing a foundation for the construction of scalable, low-cost, and high-reliability cognitive agent systems in the future.